\section{Conclusion and outlook}
\label{sec:conclusion}

We documented the Meta-Quest teleoperation stack of the
\texttt{so101\_garment} dual SO-101 rig: clutched absolute-position
retargeting, the armplane orientation construction that sidesteps the 5-DoF
over-constraint, an analytic workspace envelope with four out-of-envelope
policies, a pluggable IK layer now including a Telegrip-style split IK, and
a decoupled thumbstick wrist interface that hands wrist roll and flex to the
operator as clutched joint-space trims. Simulation results support two
concrete recommendations:
\begin{itemize}
  \item \textbf{Wrist agility.} The Telegrip split IK halves roll lag and
  reduces orientation error relative to the armplane QP at every tested
  frequency, and simultaneously improves position tracking; its higher
  command jerk argues for hardware validation (servo wear, audible
  chatter) before promotion to the pipeline default. Tuning the QP
  alternative (frame-task gain, One-Euro \(\beta\)) along the wrist-agility
  suite is the complementary path.
  \item \textbf{Out-of-envelope handling.} \texttt{project} should replace
  the legacy pass-through: bounded error, boundary sliding, seamless
  re-entry, negligible cost. \texttt{freeze} is a useful conservative
  alternative for delicate scenes; \texttt{slow} adds a tuning burden for
  marginal benefit at teleop speeds.
\end{itemize}
Next steps: executing the bimanual user study of
Section~\ref{sec:exp-userstudy}, whose subjective measures (workload,
usability, teleoperation feel) complement the objective benchmark and
arbitrate the wrist-agility/jerk trade on real operators; folding the
envelope margin into operator feedback (controller vibration); tuning the
thumbstick wrist interface (trim rates, deadzone, and direction signs) from
operator feel on the real rig, where its parameters are so far only
guesses; and extending the envelope to a self-collision-aware region for
bimanual work near the rig center.
